\chapter{Phụ lục thiết kế thay thế và độ nhạy}
\label{app:thiet-ke}

Phụ lục này phân tích các lựa chọn thiết kế và phương án thay thế. Mục đích không phải tuyên bố
phương án đã chọn là tối ưu, mà để người đọc biết kết quả có thể thay đổi thế nào nếu đổi tham
số. Những chỗ có số liệu được ghi rõ; những chỗ chỉ là phân tích lý thuyết được ghi là giả
thuyết.

\section{Tập ứng viên và bán kính}
\paragraph{Hiện tại.} Các cấp 2--4 dùng \texttt{candidates(radius=2)}, tức ô trống cách quân gần
nhất không quá 2 theo Chebyshev. Đây là cắt tỉa tiến: nó giảm hệ số rẽ nhánh nhưng có thể loại
nước tốt.
\paragraph{Phương án.} Bán kính 1 (hẹp hơn), 3 (rộng hơn), hoặc toàn bàn (không cắt).
\paragraph{Phân tích.} Với bán kính 1, số ứng viên giảm khoảng một nửa, nhưng các nước tấn công
xa (ví dụ chặn một dãy chéo ở khoảng cách 2) có thể bị bỏ. Với bán kính 3, số ứng viên tăng
mạnh, làm cây sâu hơn nhưng chậm hơn. Trong Caro, hầu hết nước chiến thuật nằm trong bán kính
2 của một quân đã có; bán kính 2 là cân bằng hợp lý. Tuy nhiên, báo cáo chưa đo ablation bán
kính, nên không thể nói bán kính 2 tốt hơn 1 hay 3 một cách định lượng.
\paragraph{Khuyến nghị.} Nếu có thời gian, chạy A9: Alpha-Beta độ sâu 4 với bán kính 1, 2, 3
trên cùng 100 khai cuộc. Đây là ablation rẻ và trả lời trực tiếp câu hỏi cắt tỉa tiến có hại
không.

\section{Trọng số hàm lượng giá}
\paragraph{Hiện tại.} \texttt{WEIGHT = (0, 1, 8, 64, 512)} cho số quân 0--4 trong cửa sổ sống.
\paragraph{Phương án.} Trọng số tuyến tính, bậc hai, hoặc học được; bỏ thành phần phòng thủ
(\texttt{no\_defense}).
\paragraph{Bằng chứng.} A7 cho thấy bỏ phòng thủ thay đổi giá trị gốc trên 9/10 thế, tức thành
phần này có tác dụng. Tuy nhiên, A7 không đo sức chơi; A7b được thiết kế để làm việc đó nhưng
chưa chạy xong. Trọng số hiện tại là lựa chọn của người thiết kế, không phải kết quả tối ưu.
\paragraph{Phân tích.} Trọng số tăng theo cấp số nhân phản ánh việc một tứ hở nguy hiểm hơn
nhiều ba hở. Nếu dùng trọng số tuyến tính, tác tử có thể đánh giá thấp đe doạ 4 quân. Nếu dùng
trọng số quá lớn, tác tử có thể bỏ qua vị trí trung tâm. Báo cáo không tối ưu trọng số; đây là
hướng mở rộng.

\section{Hệ số khám phá UCT}
\paragraph{Hiện tại.} \texttt{c\_uct = 1.4}.
\paragraph{Phương án.} 0.5 (khai thác mạnh), 1.0, 2.0 (khám phá mạnh).
\paragraph{Phân tích.} $c$ nhỏ làm MCTS tập trung vào nước đang có kết quả tốt; $c$ lớn làm nó
thử nhiều nước hơn. Với phần thưởng thưa và phương sai lớn như Caro, $c$ quá nhỏ dễ kẹt ở một
nhánh; $c$ quá lớn dễ lãng phí mô phỏng. Giá trị 1.4 là mặc định phổ biến, nhưng chưa được
tối ưu trong dự án. A3/A4 giữ cố định 1.4, nên không đo độ nhạy.
\paragraph{Khuyến nghị.} Chạy A10: MCTS 400 mô phỏng với $c \in \{0.5, 1.0, 1.4, 2.0\}$ trước
L2. Rẻ và có thể cho thấy một phần lý do MCTS yếu.

\section{Hệ số PUCT và tiên nghiệm}
\paragraph{Hiện tại.} \texttt{c\_puct = 1.5}, \texttt{leaf\_eval = value}, \texttt{batch\_leaves = 8}.
\paragraph{Phương án.} $c_{\text{puct}} \in \{0.5, 1.0, 2.5\}$; \texttt{leaf\_eval = mix} với
$\lambda=0.5$; \texttt{batch\_leaves} 1 hoặc 32.
\paragraph{Phân tích.} $c_{\text{puct}}$ điều chỉnh mức tin vào prior. Nếu prior tốt, $c$ nhỏ
giúp khai thác; nếu prior kém, $c$ lớn giúp tìm kiếm sửa sai. Trong dự án, prior đến từ SL v3
(top-1 67,8\%), không hoàn hảo, nên $c=1.5$ là lựa chọn trung dung. \texttt{leaf\_eval=mix}
trộn value với rollout, có thể giảm phương sai nhưng thêm chi phí; chưa đo. \texttt{batch\_leaves}
ảnh hưởng thông lượng, không ảnh hưởng kết quả lý thuyết nếu virtual loss đúng.
\paragraph{Khuyến nghị.} Chạy A11: L7 200 mô phỏng với $c_{\text{puct}} \in \{0.5,1.0,1.5,2.5\}$
trước L3, 100 ván mỗi mức. Đây là ablation có khả năng tạo khác biệt lớn.

\section{Dữ liệu chuyên gia}
\paragraph{Hiện tại.} \texttt{expert\_v2}: 20\,000 ván, top-k 5, temperature 8.0, depth 3.
\paragraph{Phương án.} Top-k 3 với temperature 0.5 (v1), top-k 5 với temperature 1, depth 4.
\paragraph{Bằng chứng.} SL v1 (v1 data) top-1 44,0\%; SL v2 (v2 data) top-1 65,3\%. Đây là bằng
chứng mạnh rằng đa dạng dữ liệu quan trọng hơn số epoch. Tuy nhiên, chưa tách riêng tác dụng
của top-k và temperature; cả hai thay đổi cùng lúc.
\paragraph{Phân tích.} Temperature cao làm nhãn mềm hơn, mạng học nhiều nước hơn nhưng có thể
học cả nước kém. Top-k lớn tăng đa dạng nhưng giảm chất lượng trung bình. Cân bằng tốt phụ
thuộc mục tiêu: nếu dùng policy cho PUCT, đa dạng có ích; nếu dùng policy-only, nhãn sắc có
thể tốt hơn.
\paragraph{Khuyến nghị.} Chạy A12: huấn luyện 4 mô hình với các cặp (top-k, temperature)
(3; 0,5), (3; 8), (5; 0,5) và (5; 8), 15 epoch, đấu 100 ván. Đây là ablation đắt nhưng trả lời
câu hỏi quan trọng về dữ liệu.

\section{Tự chơi và gating}
\paragraph{Hiện tại.} 200 ván/vòng, 400 mô phỏng, Dirichlet $\alpha=0.3$, $\varepsilon=0.25$,
temperature 1 cho 8 nước đầu, gating 100 ván ngưỡng 55\%.
\paragraph{Phương án.} Nhiều ván hơn, ít mô phỏng hơn, Dirichlet khác, ngưỡng 50\% hoặc 60\%.
\paragraph{Bằng chứng.} L8 chỉ duyệt 6/50 vòng. Ngưỡng 55\% trên 100 ván có lực lượng thấp:
với chênh lệch thật 5 điểm, xác suất vượt 55\% không cao. Điều này có thể làm quá trình bảo
thủ, giữ bản cũ dù ứng viên tốt hơn. Ngược lại, ngưỡng thấp có thể chấp nhận ứng viên kém.
\paragraph{Phân tích.} Gating là một kiểm định thống kê có sai số. Với 100 ván, khoảng Wilson
của 55\% là khoảng 45--65\%, nên một ứng viên thật 55\% có thể chỉ đạt 45\%. Đây là lý do nên
dùng nhiều ván hơn hoặc tiêu chí mềm hơn (ví dụ trung bình trượt).
\paragraph{Khuyến nghị.} Nếu chạy lại L8, tăng gating lên 200 ván và dùng ngưỡng 52--55\%, hoặc
đánh giá trên tập khai cuộc giữ riêng. Chi phí tăng không lớn so với tự đấu.

\section{Nhiệt độ và chọn nước khi thi đấu}
\paragraph{Hiện tại.} Policy-only dùng argmax; PUCT dùng nước nhiều lượt thăm nhất; tự chơi
dùng temperature 1 cho 8 nước đầu.
\paragraph{Phương án.} Lấy mẫu theo $\pi^{1/\tau}$ khi thi đấu, hoặc temperature thấp.
\paragraph{Phân tích.} Argmax tối đa hoá xác suất nhưng có thể bỏ qua nước có xác suất thấp mà
thắng. Lấy mẫu tăng đa dạng nhưng giảm sức mạnh trung bình. Trong thi đấu, argmax thường tốt
hơn; trong tự chơi, lấy mẫu cần thiết để khám phá. Đây là lý do cấu hình tách hai chế độ.

\section{Kích thước bàn và luật}
\paragraph{Hiện tại.} Bàn 11×11, luật tự do.
\paragraph{Phương án.} Bàn 9×9, 15×15, luật chặn hai đầu, Renju.
\paragraph{Phân tích.} Bàn lớn hơn làm hệ số rẽ nhánh và độ dài ván tăng; bàn nhỏ hơn làm lợi
thế đi trước mạnh hơn. Luật chặn hai đầu làm mẫu đe doạ phức tạp hơn, có thể làm Alpha-Beta
mạnh hơn tương đối vì tri thức viết tay quan trọng hơn. Renju có luật cấm, cần mô hình hoá
trạng thái phức tạp hơn. Báo cáo không thử các biến thể này; kết quả chỉ áp dụng cho cấu hình
đã dùng.

\section{Bảng tóm tắt}
\begin{small}
\begin{longtable}{@{}p{.22\textwidth} p{.18\textwidth} p{.14\textwidth} p{.34\textwidth}@{}}
\toprule
\textbf{Tham số} & \textbf{Hiện tại} & \textbf{Đã đo?} & \textbf{Nhận xét} \\
\midrule
Bán kính ứng viên & 2 & Không & Nên chạy A9 \\
Trọng số lượng giá & (1,8,64,512) & A7 một phần & Cần A7b \\
$c_{\text{uct}}$ & 1.4 & Không & Nên chạy A10 \\
$c_{\text{puct}}$ & 1.5 & Không & Nên chạy A11 \\
Dữ liệu expert & top-5, T=8 & Có (SL v1/v2) & Đa dạng giúp top-1 \\
Dirichlet & 0.3/0.25 & Không & Ảnh hưởng đa dạng tự chơi \\
Gating & 100 ván, 55\% & Có (L8) & Có thể quá bảo thủ \\
Nhiệt độ thi đấu & argmax & Không & Hợp lý cho thi đấu \\
Bàn/luật & 11×11 tự do & Không & Ngoài phạm vi \\
\bottomrule
\end{longtable}
\end{small}